\chapter{Parameter Estimation and Model Fitting}
\label{ch:param-estimation}

\epigraph{All models are wrong, but some are less wrong than others once you fit them.}{}

\section{System Identification for Multi-Body Models}

\begin{lstlisting}[language=Python, caption={Inertial parameter estimation from motion data.}]
import numpy as np
from scipy.optimize import minimize

def estimate_inertial_params(
    q_data: np.ndarray,
    qd_data: np.ndarray,
    qdd_data: np.ndarray,
    tau_data: np.ndarray,
    n_params: int
) -> np.ndarray:
    """Estimate inertial parameters using least-squares.

    The equations of motion are LINEAR in inertial parameters:
    tau = Y(q, qd, qdd) * pi
    where Y is the regressor matrix and pi is the parameter vector.

    Parameters
    ----------
    q_data : (N, n) joint positions
    qd_data : (N, n) joint velocities
    qdd_data : (N, n) joint accelerations
    tau_data : (N, n) joint torques
    n_params : int number of inertial parameters

    Returns
    -------
    np.ndarray estimated parameters
    """
    N, n = q_data.shape

    # Build regressor matrix (simplified for 1-DOF)
    Y = np.zeros((N * n, n_params))
    tau_vec = tau_data.reshape(-1)

    for k in range(N):
        for j in range(n):
            row = k * n + j
            Y[row, 0] = qdd_data[k, j]  # Inertia
            Y[row, 1] = np.sin(q_data[k, j])  # Gravity

    # Least squares solution
    params, _, _, _ = np.linalg.lstsq(Y, tau_vec, rcond=None)
    return params
\end{lstlisting}

\section{Muscle Parameter Calibration}

Using the Hill model from Volume~III with optimization to match experimental data.

\section*{Exercises}
\begin{enumerate}
  \item Estimate the mass and length of a pendulum from simulated motion data.
  \item Add measurement noise and study estimation robustness.
  \item Implement Bayesian estimation with prior knowledge from literature.
\end{enumerate}
